\section{Конструкорская часть}

\subsection{Требования к программному обеспечению}
Программное обеспечение должно обеспечивать следующие функциональные возможности:
\begin{enumerate}
    \item добавление геометрического объекта из списка: сфера, цилиндр, параллелепипед,
    трехугольная и четырехугольная пирамида, конус, цилиндр;
    \item добавление нескольких источников света;
    \item добавление нескольких камер, с возможностью переключения между ними;
    \item изменения материала объекта.
\end{enumerate}

\subsection{Модули программного обеспечения}

Структура разрабатываем программы описана на рисунке~\ref{fig:struct-diag}.
\begin{figure}[H]
	\centering
	\includegraphics[width=1\textwidth]{diag/struct-scheme.pdf}
	\caption{Структурная схема программы}
	\label{fig:struct-diag}
\end{figure}

Разрабатываемую программу можно поделить на следующие модули:
\begin{enumerate}
    \item интерфейс --- графическая оболочка для взаимодействия с пользователем.
    \item Модуль сцены --- предоставляет методы для редактирования объектов сцены.
    \item Модуль генерации изображения --- генерирует итоговое изображение.
\end{enumerate}

\subsection{Разработка алгоритмов}
\subsubsection{Алгоритм трассировки пути}

Процесс генерации изображения выполняется для заданной области, что обеспечивает возможность параллельной обработки.
Для каждого пикселя в области испускается заданное количество лучей, и результаты усредняются для снижения шума методом Монте-Карло.
Схема реализации алгоритма генерации части изображения представлена на рисунке~\ref{fig:risunok_1} 
\begin{figure}[H]
	\centering
	\includegraphics[width=0.75\textwidth]{diag/2-1.pdf}
	\caption{Схема реализации алгоритма генерации части изображения}
	\label{fig:risunok_1}
\end{figure}

Цвет пикселя с координатами~$(i, j)$ вычисляется как среднее значение цвета по всем лучам:
\begin{equation}
C_{i,j} = \frac{1}{N_s} \sum_{k=1}^{N_s} L(r_k),
\end{equation}
\noindent где $N_s$~--- количество лучей на пиксель, $L(r_k)$~--- яркость, возвращаемая функцией трассировки для луча~$r_k$.

\subsubsection{Алгоритм генерации луча}
Для каждого пикселя генерируется луч, проходящий через случайную точку внутри него,
что реализует сглаживание. Координаты точки на плоскости изображения вычисляются как~\cite{Shirley2025RTW2}:
\begin{equation}
p_{\text{sample}} = p_{00} + (i + \xi_x) \Delta u + (j + \xi_y) \Delta v,
\end{equation}

\noindent где $p_{00}$~--- позиция центра пикселя~$(0, 0)$, $\Delta u$ и~$\Delta v$~--- векторы, определяющие размер пикселя в горизонтальном и вертикальном направлениях, $\xi_x, \xi_y \sim U(0, 1)$~--- случайные смещения для выборки Монте-Карло.

Начало луча определяется с учётом эффекта глубины резкости.
Для камер, где угол размытия больше нуля начало луча случайным образом выбирается на диске размытия.

Направление луча вычисляется как:
\begin{equation}
d = p_{\text{sample}} - ray\_origin,
\end{equation}
где ray\_origin это точка начала луча.

Схема алгоритма генерации луча представлена на рисунке~\ref{fig:risunok_2}
\begin{figure}[H]
	\centering
	\includegraphics[width=0.75\textwidth]{diag/2-2.pdf}
	\caption{Схема алгоритма генерации луча}
	\label{fig:risunok_2}
\end{figure}


\subsubsection{Рекурсивная трассировка}

Вычисление цвета луча выполняется рекурсивно с ограничением глубины рекурсии для предотвращения бесконечных отражений. Алгоритм реализует следующую логику:

\begin{enumerate}
    \item проверка глубина рекурсии. Если глубина рекурсии превышена, возвращается черный цвет;
    \item проверка пересечения луча со всеми объектами сцены. Если пересечений нет, возвращается черный цвет;
    \item вычисления излучения материала в точке пересечения. Формула~(\ref{eq:emit});
    \item рассеивание луча. Материал определет, происходит ли рассеивание луча, вычисляет новое направление, и коэффициент затухания;
    \item рекурсивный вызов функции. Для рассеяного луча рекурсивно вычисляется цвет с уменьшенной глубиной. Формула~(\ref{eq:req-scatter}).
\end{enumerate}

\begin{equation}
    \label{eq:emit}
C_{\text{emit}} = L_e(p),
\end{equation}
\noindent где $p$~--- координаты точки пересечения.


\begin{equation}
    \label{eq:req-scatter}
C_{\text{scatter}} = a \odot L(r', d - 1),
\end{equation}
\noindent где $a$ --- цвет луча после столкновения, $\odot$~--- покомпонентное произведение векторов (цветов), $r'$~--- рассеянный луч с началом в точке пересечения, $d$ --- количество переотражений луча.

Финальный цвет луча определяется как сумма излучения и рассеянного света:
\begin{equation}
L(r, d) = C_{\text{emit}} + C_{\text{scatter}}.
\end{equation}

Эта формула является дискретной аппроксимацией уравнения рендеринга, решаемого методом Монте-Карло.

Схема алгоритма рекурсивной трассировки пути представлена на рисунке~\ref{fig:risunok_3}
\begin{figure}[H]
	\centering
	\includegraphics[width=0.65\textwidth]{diag/2-3.pdf}
	\caption{Схема алгоритма рекурсивной трассировки пути}
	\label{fig:risunok_3}
\end{figure}

Рекурсивная структура алгоритма естественным образом моделирует глобальное освещение:
луч может многократно отражаться от поверхностей, собирая непрямой свет, что создаёт реалистичные эффекты цветового переотражения и мягких теней.

\subsection{Алгоритм визуализации тумана}

Взаимодействие луча с объёмным туманом моделируется стохастически с учётом переменной плотности среды, что обеспечивает реалистичное поведение неоднородного тумана~\cite{volumetric-fog}.

\textbf{Проверка пересечения с границами.} Сначала определяется интервал~$[t_0, t_1]$, на котором луч находится внутри ограничивающего объёма тумана:

\begin{equation}
t_0 = \max(t_{\text{entry}}, t_{\min}), \quad t_1 = \min(t_{\text{exit}}, t_{\max}),
\end{equation}

\noindent где $t_{\text{entry}}$ и~$t_{\text{exit}}$~--- параметры входа и выхода луча из ограничивающего объёма, $[t_{\min}, t_{\max}]$~--- допустимый интервал трассировки. Если~$t_0 \geq t_1$, луч не проходит через объём тумана.

\textbf{Вычисление расстояния до взаимодействия.} Расстояние, на котором происходит взаимодействие фотона с частицами тумана, определяется методом инверсной функции распределения для экспоненциального закона~\cite{pbr-book}:

\begin{equation}
d_{\text{hit}} = -\frac{\ln(\xi)}{\rho_0},
\end{equation}
\noindent где $\xi \sim U(0, 1)$~--- случайная величина, $\rho_0$~--- базовая плотность тумана. Параметр луча, соответствующий точке взаимодействия:

\begin{equation}
t = t_0 + \frac{d_{\text{hit}}}{|d|},
\end{equation}

\noindent где $|d|$~--- длина вектора направления луча. Если~$d_{\text{hit}} > (t_1 - t_0) \cdot |d|$, взаимодействие не происходит в пределах объёма.

\textbf{Вычисление локальной плотности.} В точке взаимодействия~$p = o + td$ вычисляется фактическая плотность с учётом неоднородности:

\begin{equation}
\rho(p) = \rho_0 \cdot f_{\text{noise}}(p) \cdot f_{\text{height}}(p),
\end{equation}

\noindent где $f_{\text{noise}}(p)$~--- функция шума Перлина, создающая турбулентность:

\begin{equation}
f_{\text{noise}}(p) = 0.2 + 0.8 \cdot |\text{noise}(s \cdot p)|,
\end{equation}

\noindent $s$~--- масштабный коэффициент, а~$f_{\text{height}}(p)$~--- экспоненциальный множитель спада по высоте для приземного тумана:

\begin{equation}
f_{\text{height}}(p) = \exp(-k \cdot \max(y, -10)),
\end{equation}

\noindent где $k > 0$~--- коэффициент спада по высоте, $y = p_y$~--- вертикальная координата.

\textbf{Стохастическое отсеивание.} Поскольку базовое расстояние~$d_{\text{hit}}$ вычислено для постоянной плотности~$\rho_0$, а фактическая плотность~$\rho(p)$ может быть меньше, применяется стохастическое отсеивание~\cite{pbr-book}:

\begin{equation}
P_{\text{accept}} = \min\left(\frac{\rho(p)}{\rho_0}, 1\right) = \min(\rho(p), 1).
\end{equation}

Генерируется случайная величина~$\xi' \sim U(0, 1)$; если~$\xi' > P_{\text{accept}}$ или $\rho(p) < 10^{-5}$, взаимодействие отвергается и поиск продолжается.

Схема алгоритма генерации тумана представлена на рисунке~\ref{fig:risunok_4}
\begin{figure}[H]
	\centering
	\includegraphics[width=0.65\textwidth]{diag/2-4.pdf}
	\caption{Схема алгоритма генерации тумана}
	\label{fig:risunok_4}
\end{figure}

\subsection*{Вывод}

В данной части были разработаны алгоритмы системы визуализации тумана методом трассировки путей.
Описаны алгоритмы генерации изображения, формирования лучей и рекурсивной трассировки, представлена схема спроектированного программного обеспечения.


\clearpage